\documentclass[twocolumn]{iopjournal}

\usepackage{amsmath,amssymb,amsfonts}
\usepackage{amsthm}

\fancyhfoffset[L]{0mm}
\setlength{\marginparwidth}{0pt}
\setlength{\marginparsep}{0pt}

% ==========================================
% ZENODO PREPRINT HEADERS (ACTIVE)
% ==========================================
\renewcommand{\articletype}[1]{%
  {\vspace*{-8mm}\noindent \Large \sf Preprint}\par
  \vspace*{6mm}\noindent{\scriptsize \sf{\bfseries \MakeUppercase{#1}}}}

\fancyhead[L]{}
\fancyhead[R]{}

\hypersetup{pdftitle={The Gram Matrix of the Tetrad Equation}, pdfauthor={Stephen Nelson}, pdfsubject={Gravity Field Dynamics (GFD)}, pdfkeywords={tetrad, Gram matrix, spatial parity, stellated octahedron, torsion, proof}}

\newtheorem{lemma}{Lemma}
\newtheorem{theorem}{Theorem}
\newtheorem{proposition}{Proposition}
\newtheorem{corollary}{Corollary}
\theoremstyle{definition}
\newtheorem{definition}{Definition}
\theoremstyle{remark}
\newtheorem*{remark}{Remark}

\begin{document}

\articletype{Supplementary Material}

\title{The Gram Matrix of the Tetrad Equation}

\author{Stephen Nelson}

\affil{Field Dynamics Research}

\ead{stephen@fielddynamics.org}

\begin{abstract}
The tetrad equation $g = e^T\eta\,e$ is even in the spatial frame,
so its spatial realisation $M$ comes with a parity partner $-M$. We
study the Gram matrices of the pair. Self-coupling is root-blind and
cross-coupling reverses every sign; since the cross-coupling is also
unchanged by exchanging the roots, Gram data resolve whether two
factors come from the same root or from opposite roots, never which
root is which. For the tetrahedral frame, bordering $M$ with a row of
ones gives a Hadamard matrix, which yields $M^TM = 4I - J$, the
spectral split $I = \tfrac14 G + \tfrac14 J$ and the rows of $M$ as
eigenvectors, at once. The Gram matrix fixes $M$ up to a unique
element of $O(3)$, and that element is $-I$ for the two roots: the
roots are told apart by orientation alone, which flips because three
is odd. The two roots are the parity classes $xyz = \pm1$ of the cube,
and the symmetry group of their compound, the stellated octahedron,
is that of one tetrahedron times $\{\pm I\}$. On the connection,
exchange of the lower indices splits the 64 components into 40 and
24; the torsion-free condition is the projection onto the 40, and for
a metric connection the discarded 24, the torsion, are the only place
a connection can tell the roots apart. A connection carried
covariantly with the frame does not; one that keeps its spin
connection in general does, through its torsion. Every step is exact and
machine-checked.
\end{abstract}

\keywords{tetrad, Gram matrix, spatial parity, Hadamard matrix,
stellated octahedron, torsion, proof}

% ======================================================================

\section*{Setting}

A Lorentz tetrad $e$ satisfies the tetrad equation
\begin{equation}
  g = e^T\eta\,e, \qquad \eta = \mathrm{diag}(-1,+1,+1,+1).
  \label{eq:tetrad}
\end{equation}
Spatial parity $P = \mathrm{diag}(1,-1,-1,-1)$ acts on the frame
index, $e \mapsto Pe$. Since $P^T\eta P = \eta$, the metric is
unchanged, $(Pe)^T\eta(Pe) = g$, while $\det(Pe) = -\det e$: the
equation has two roots, carried into each other by $P$, and the
metric sees only $(\det e)^2$~\cite{NelsonScheme}.

Following~\cite{NelsonLetter1,NelsonScheme}, the frame is realised
in an observer's space by four directions $v_1,\dots,v_4$, the
columns of a matrix $M \in \mathbb{R}^{3\times4}$, each a fixed linear
combination of the spatial legs $e^1, e^2, e^3$. Every pair of
directions bears the same relation and every direction has the same
magnitude; these premises force the Gram matrix $M^TM$ and the
regular tetrahedron, unique up to $O(3)$~\cite{NelsonForced}. We take
that result as given, scale to $|v_a|^2 = 3$ (every statement below
is invariant under a common rescaling), and use the representative
\begin{equation}
  M_1 = \begin{pmatrix}
    +1 & +1 & -1 & -1 \\
    +1 & -1 & +1 & -1 \\
    +1 & -1 & -1 & +1
  \end{pmatrix}.
  \label{eq:M1}
\end{equation}
Since each $v_a$ is linear in the spatial legs, $P$ sends every
$v_a \mapsto -v_a$, so the second root is realised by $M_2 = -M_1$.
Reversing every leg, $e \mapsto -e = PTe$, also reverses the time
leg; as in~\cite{NelsonScheme}, the identification of root with
tetrahedron is for future-pointing frames. The legs are labelled
throughout: $v_a$ and $-v_a$ belong to the same leg. We write $I$ for the identity,
$J$ for the $4\times4$ all-ones matrix and $\mathbf{1} = (1,1,1,1)^T$.

% ======================================================================

\section*{Part I: The three Gram products}

\begin{definition}
For realisations $U, W$ the Gram product is
$G_{UW} = U^T W$, with entries $v_a \cdot w_b$. From the parity pair
we form $G_{M_1M_1}$, $G_{M_2M_2}$, $G_{M_1M_2}$ and
$G_{M_2M_1} = G_{M_1M_2}^T$. We write $G = G_{M_1M_1}$.
\end{definition}

\begin{lemma}[Root-blindness]
\label{lem:blind}
$G_{M_2M_2} = G_{M_1M_1}$.
\end{lemma}

\begin{proof}
$(-M_1)^T(-M_1) = M_1^TM_1$.
\end{proof}

\begin{lemma}[Cross-coupling]
\label{lem:cross}
$G_{M_1M_2} = G_{M_2M_1} = -G$. Hence the Gram products take the
value $G$ on a same-root pair and $-G$ on a cross-root pair, and each
value is unchanged when the two roots are exchanged.
\end{lemma}

\begin{proof}
$M_1^T(-M_1) = -G = (-M_1)^TM_1$. Exchanging the roots sends
$G_{M_1M_1} \leftrightarrow G_{M_2M_2}$ and
$G_{M_1M_2} \leftrightarrow G_{M_2M_1}$, which are equal pairs.
\end{proof}

\noindent
So the cross-coupling is \emph{root-resolving} in a precise sense: it
resolves the \emph{relative} root of two factors. No Gram product
says which root a single frame is; Proposition~\ref{prop:orient}
gives the quantity that does.

\begin{lemma}[Explicit Gram matrix]
\label{lem:explicit}
Let $\hat H = \bigl(\begin{smallmatrix}\mathbf{1}^T \\ M_1
\end{smallmatrix}\bigr)$. Then $\hat H\hat H^T = 4I$, hence
$\hat H^T\hat H = 4I$, $M_1\mathbf{1} = 0$ and
\begin{equation}
  G = M_1^TM_1 = 4I - J, \qquad M_1M_1^T = 4I_3,
  \label{eq:G}
\end{equation}
\end{lemma}

\begin{proof}
The rows of $M_1$ sum to zero and are mutually orthogonal with
squared norm 4, by inspection of~(\ref{eq:M1}); with the row
$\mathbf{1}^T$, of squared norm 4, this is $\hat H\hat H^T = 4I$.
A square matrix with $\hat H\hat H^T = 4I$ has
$\hat H^{-1} = \tfrac14\hat H^T$, so $\hat H^T\hat H = 4I$. Its
left side is $\mathbf{1}\mathbf{1}^T + M_1^TM_1 = J + G$.
\end{proof}

\noindent
$\hat H$ is the $4\times4$ Hadamard matrix~\cite{NelsonHadamard}.
Entrywise, $G$ has diagonal $3$ and off-diagonal $-1$; with
Lemmas~\ref{lem:blind} and~\ref{lem:cross} all three products are
determined:
\begin{equation}
  G_{M_1M_1} = G_{M_2M_2} = 4I - J, \qquad G_{M_1M_2} = J - 4I.
  \label{eq:three}
\end{equation}

\begin{lemma}[Eigenstructure]
\label{lem:eigen}
$\tfrac14G$ and $\tfrac14J$ are complementary orthogonal projections,
$I = \tfrac14G + \tfrac14J$. Thus $G$ has eigenvalue $0$ on
$\mathbf{1}$ (closure) and eigenvalue $4$ with multiplicity $3$ on
$\mathbf{1}^\perp$, which is spanned by the rows of $M_1$:
$G M_1^T = 4M_1^T$. The cross-coupling has eigenvalues $0$ and $-4$
on the same spaces.
\end{lemma}

\begin{proof}
$J^2 = 4J$ gives $(\tfrac14J)^2 = \tfrac14J$, and
$\tfrac14G = I - \tfrac14J$ is the complementary projection.
$GM_1^T = M_1^T(M_1M_1^T) = 4M_1^T$, and the three rows are
independent by $M_1M_1^T = 4I_3$.
\end{proof}

\begin{lemma}[Invariance and rigidity]
\label{lem:orth}
For $Q \in O(3)$ the three products are unchanged by
$(M_1,M_2) \mapsto (QM_1, QM_2)$. Conversely, if $W$ has
$W^TW = G$, then $W = QM_1$ for a unique $Q \in O(3)$, namely
$Q = \tfrac14 WM_1^T$. For $W = M_2$, $Q = -I$.
\end{lemma}

\begin{proof}
$(QU)^T(QW) = U^TQ^TQW = U^TW$. The converse is the rigidity lemma
of~\cite{NelsonForced}, since $G$ has rank 3; multiplying
$W = QM_1$ on the right by $M_1^T$ and using $M_1M_1^T = 4I_3$ gives
$Q$. For $W = -M_1$ this is $-I$.
\end{proof}

\noindent
So, for realisations of rank 3 with labelled legs, the Gram matrix
is a complete invariant under $O(3)$, and the two roots are one
$O(3)$ orbit: they are related by
$-I$, which has determinant $-1$ and lies outside $SO(3)$.

\begin{proposition}[Orientation]
\label{prop:orient}
For each triple of legs $a<b<c$, the orientation
$\omega_{abc}(M) = \det(v_a, v_b, v_c)$ is nonzero on $M_1$,
invariant under $SO(3)$, and satisfies
$\omega_{abc}(M_2) = -\omega_{abc}(M_1)$.
\end{proposition}

\begin{proof}
Any three columns of $M_1$ are independent, since $\mathbf{1}$ spans
the null space of $M_1$. $\det(Qv_a,Qv_b,Qv_c) = \det Q\,\omega_{abc}$,
and $\det(-I_3) = (-1)^3 = -1$.
\end{proof}

\noindent
This is the spatial counterpart of $\det(Pe) = -\det e$: the roots
differ only in orientation, and orientation flips because $D = 3$ is
odd.

% ======================================================================

\section*{Part II: Geometric identification}

\begin{theorem}[Tetrahedra in the cube]
\label{thm:cube}
The columns of $M_1$ are the four corners of the cube $\{\pm1\}^3$
with $xyz = +1$, and those of $M_2$ the four with $xyz = -1$. Each
set is a regular tetrahedron centred at the origin, with every edge
of squared length $8$. The two sets are disjoint and together are the
eight corners of the cube; their compound is the stellated
octahedron $S = [\,M_1 \mid M_2\,]$~\cite{Kepler1619,Coxeter1973}.
\end{theorem}

\begin{proof}
Each column of~(\ref{eq:M1}) has an even number of entries $-1$, so
$xyz = +1$; negating three coordinates multiplies $xyz$ by
$(-1)^3$. The corners with $xyz = +1$ are four in number, so they are
exactly the columns of $M_1$, and the rest are those of $M_2$. By
Lemma~\ref{lem:explicit},
$|v_a - v_b|^2 = 3 + 3 - 2(-1) = 8$ for $a \neq b$, and
$\sum_a v_a = M_1\mathbf{1} = 0$.
\end{proof}

\noindent
Following~\cite{NelsonLetter1,NelsonScheme}, the first tetrahedron is
TetA (the identity component under spatial parity) and the second
TetB (the parity component).

\begin{theorem}[Symmetry and non-isomorphism]
\label{thm:sym}
The orthogonal maps permuting the corners of TetA form a group of
order 24, acting on the four legs as the whole of $S_4$. Those
permuting the corners of $S$ form the group
\begin{equation}
  \mathrm{Sym}(S) = \mathrm{Sym}(\mathrm{TetA}) \times \{I, -I\},
  \label{eq:sym}
\end{equation}
of order 48, and every element outside $\mathrm{Sym}(\mathrm{TetA})$
carries TetA onto TetB.
\end{theorem}

\begin{proof}
A map in $\mathrm{Sym}(\mathrm{TetA})$ is fixed by the permutation it
induces, as the corners span $\mathbb{R}^3$. Every permutation
occurs: by Lemma~\ref{lem:orth}, $M_1\Pi$ has Gram matrix
$\Pi^TG\Pi = G$ for any permutation matrix $\Pi$, so
$M_1\Pi = QM_1$ for some $Q \in O(3)$. Hence the order is 24.
Two corners of the cube are at squared distance $8$ exactly when
they differ in two coordinates, that is, when they have the same
$xyz$. A map $L \in \mathrm{Sym}(S)$ is an isometry, so it carries
TetA, four corners at mutual squared distance $8$, onto four such
corners, all with one value of $xyz$: onto TetA or onto TetB, and
then TetB $= -$TetA onto the other. Those of the first kind form
$\mathrm{Sym}(\mathrm{TetA})$. If $L$ is of the second kind, $-L$ is
of the first, since $-I$ exchanges TetA and TetB; and $-I$ commutes
with every map, which gives~(\ref{eq:sym}). The verification script
confirms the count exhaustively over all $8!$ relabellings of the
corners.
\end{proof}

\noindent
The extra factor is exactly parity: the compound has twice the
symmetry of either root, and the doubling is the element that
exchanges them. Half of the exchanging elements are rotations (for
example a quarter turn about a coordinate axis), so TetA and TetB are
congruent as unlabelled point sets; as labelled frames they are not,
by Proposition~\ref{prop:orient}. This is the caution
of~\cite{NelsonScheme}: the roots differ in the order of the legs,
not in shape.

% ======================================================================

\section*{Part III: The index grid}

The affine connection $\Gamma^\lambda{}_{\mu\nu}$ has $K^3 = 64$
components. Fixing $\lambda$ gives a $4\times4$ \emph{slice} in
$(\mu,\nu)$, with $4$ diagonal entries, $6$ in the strict upper
triangle and $6$ in the strict lower triangle. Over the four slices,
\begin{equation}
  \underbrace{4\times6}_{24} + \underbrace{4\times4}_{16}
  + \underbrace{4\times6}_{24} = 64 .
  \label{eq:grid}
\end{equation}
In the framework's reading of the grid the upper triangle is TetA,
the lower triangle TetB and the diagonal shared, and exchanging the
triangles is the TetB step~\cite{NelsonGramOrig,NelsonChiral};
Letter~2 calls the exchange-even sector the identity sector and the
exchange-odd sector the parity sector~\cite{NelsonLetter2}.

\begin{lemma}[The exchange split]
\label{lem:tau}
Let $\tau$ exchange the lower indices,
$(\tau\Gamma)^\lambda{}_{\mu\nu} = \Gamma^\lambda{}_{\nu\mu}$. Then
$\tau^2 = 1$, $\tau$ fixes the diagonal and exchanges the two
triangles, and $\Pi_+ = \tfrac12(1+\tau)$ is a projection of rank
$40$ whose kernel, the $\tau$-odd part, has dimension $24$ and is the
torsion $T^\lambda{}_{\mu\nu} = \Gamma^\lambda{}_{\mu\nu}
- \Gamma^\lambda{}_{\nu\mu}$.
\end{lemma}

\begin{proof}
On a mirror pair with upper entry $u$ and lower entry $l$, $\tau$
swaps $u$ and $l$; the even combination $\tfrac12(u+l)$ is kept and
the odd one $u - l$, the torsion, is removed. Each slice has $4 + 6$
even and $6$ odd directions; times four slices, $40$ and $24$.
\end{proof}

\noindent
Written in the coordinates of~(\ref{eq:grid}), the torsion-free
condition is $l = u$: it keeps the diagonal and one triangle and
makes the other a copy, the form in which it reads as keeping shared
and TetA and discarding TetB. Independently of which triangle is
written down, what is kept is the exchange-even combination
$\tfrac12(u+l)$, which does not change when the triangles are
exchanged, and what is discarded is the exchange-odd difference
$u - l$.

Two cautions. First, the triangles depend on the order in which the
coordinates are listed: relabelling them in reverse carries every
upper entry into the lower triangle. The exchange $\tau$ commutes
with every relabelling, and the torsion is a tensor (the
inhomogeneous term in the transformation of $\Gamma$ is symmetric in
the lower indices), so the split $40 + 24$ and the condition $T = 0$
are coordinate-independent, while the split $24 + 16 + 24$ is a
property of a chosen ordered basis. Second, $\tau$ acts on the index
slots of the connection and $P$ on the frame; they are different
operations, and the only link between them proved here is
Theorem~\ref{thm:proj}(iii).

% ======================================================================

\section*{Part IV: The torsion-free projection}

\begin{theorem}[Torsion-free projection]
\label{thm:proj}
\begin{enumerate}
\item[(i)] The torsion-free condition is the projection $\Pi_+$ of
  the 64 components onto 40, with the 24 torsion components as
  kernel. It is not injective, so it is not an isomorphism:
  \begin{equation}
    \boxed{16 + 24 = 40 \;\neq\; 64 = 24 + 16 + 24.}
    \label{eq:noniso}
  \end{equation}
\item[(ii)] For any metric $g$, a connection compatible with $g$ is
  determined by $g$ and its torsion: the differences of two
  $g$-compatible connections form a 24-dimensional space, mapped
  bijectively onto the 24-dimensional space of torsions.
\item[(iii)] Hence, if each root is given a $g$-compatible connection,
  the two connections differ only through their torsion. A
  torsion-free metric connection is the Levi-Civita connection of
  $g$, which by~(\ref{eq:tetrad}) is the same on both roots.
\end{enumerate}
\end{theorem}

\begin{proof}
(i) is Lemma~\ref{lem:tau}. (ii) All statements are pointwise. The difference $\Delta$ of two
$g$-compatible connections satisfies
$\Delta_{\lambda\mu\nu} + \Delta_{\nu\mu\lambda} = 0$, with the first
index lowered by $g$: antisymmetry in two of three indices, $4\times6
= 24$ dimensions. If its torsion vanishes, $\Delta$ is also symmetric
in the last two, and a tensor symmetric in one pair of slots and
antisymmetric in another vanishes; so $\Delta \mapsto T(\Delta)$ is
injective between spaces of equal dimension, and bijective. Its
inverse is the contorsion~\cite{Hehl1976}. (iii) The metric is the
same on both roots (Setting), so both connections are the Levi-Civita
connection of $g$ plus the contorsion of their torsion, which
vanishes in the torsion-free case. The verification script checks (ii)
exactly for $\eta$ and for a generic Lorentzian metric.
\end{proof}

\begin{proposition}[How a connection follows the root]
\label{prop:follow}
Let $\omega_\mu$ be a Lorentz connection, $\eta\omega_\mu$
antisymmetric, and
$\Gamma^\lambda{}_{\mu\nu}
= (e^{-1})^\lambda{}_a\bigl(\partial_\mu e^a{}_\nu
+ \omega^a{}_{\mu b}\,e^b{}_\nu\bigr)$, which is compatible with
$g = e^T\eta e$.
\begin{enumerate}
\item[(a)] If $\omega$ is carried along with the frame,
  $(e,\omega) \mapsto (Pe, P\omega P)$, then $P\omega P$ is again a
  Lorentz connection and $\Gamma$, with its torsion, is unchanged.
\item[(b)] If both roots carry the same $\omega$, the two connections
  are compatible with the same $g$ and in general differ; by
  Theorem~\ref{thm:proj}(ii) they differ exactly by a contorsion,
  with nonzero torsion difference.
\end{enumerate}
\end{proposition}

\begin{proof}
(a) $P^T\eta P = \eta$ keeps $\eta P\omega P$ antisymmetric, and
$(Pe)^{-1}(P\partial e + P\omega P\,Pe) = e^{-1}(\partial e + \omega e)$
since $P^2 = 1$. (b) Both are of the stated form with the same $g$;
the script exhibits a rational example with nonzero torsion
difference.
\end{proof}

\begin{remark}
Parts I and IV say the same thing at two levels. The metric, and
every Gram product, cannot tell the roots apart; neither can any
torsion-free metric connection, since it is built from the metric.
At the connection level the only room for the distinction is the 24
torsion components that general relativity sets to
zero~\cite{LeviCivita1917} and Einstein--Cartan theory
keeps~\cite{Cartan1922,Hehl1976}. GFD reads the torsion as where
TetA and TetB are told apart~\cite{NelsonLetter2,NelsonTorsion}. The
theorem proves that the torsion is the only place available: the
torsion is the \emph{necessary} carrier of the distinction, as
Letter~2 states. Proposition~\ref{prop:follow} shows it is not
automatically a carrier: a connection carried covariantly with the
frame is identical on both roots, torsion included, and the roots
are distinguished at the connection level only when their
connections are not so related. Which assignment the field makes is
a physical input, not a consequence of the tetrad equation.
\end{remark}

% ======================================================================

\section*{Contributions and prior work}

The ingredients are classical: rigidity of Gram matrices under the
orthogonal group, the $4\times4$ Hadamard
matrix~\cite{Sylvester1867}, the stellated octahedron and its
symmetry $T_d \times \{\pm I\} = O_h$~\cite{Coxeter1973}, the
fundamental theorem of Riemannian geometry~\cite{LeviCivita1917},
and the decomposition of a metric connection into Levi-Civita
connection and contorsion~\cite{Hehl1976}. What this note adds is
their assembly on the two roots of the tetrad equation: the
distinction between relative and absolute root, with orientation as
the absolute invariant; the compound's symmetry as one root's
doubled by the exchange; and the exact form of the statement that
torsion is the only place a metric connection can tell the roots
apart, together with the condition under which it does.

% ======================================================================

\section*{Summary}

\vspace{4pt}
\noindent
{\footnotesize
\begin{tabular*}{\columnwidth}{@{\extracolsep{\fill}}lll@{}}
\hline
Result & Proves & Requires \\
\hline
Lemma 1 & $G_{M_2M_2} = G_{M_1M_1}$ & $M_2 = -M_1$ \\
Lemma 2 & $G_{M_1M_2} = G_{M_2M_1} = -G$ & idem \\
Lemma 3 & $G = 4I - J$ & Hadamard \\
Lemma 4 & $I = \tfrac14G + \tfrac14J$ & Lemma 3 \\
Lemma 5 & Invariance, rigidity & \cite{NelsonForced} \\
Prop.\ 1 & Orientation flips & $D = 3$ odd \\
Thm.\ 1 & $xyz = \pm1$; cube & Lemma 3 \\
Thm.\ 2 & Symmetry $\times\{\pm I\}$ & Lemma 5 \\
Lemma 6 & $64 = 40 + 24$ under $\tau$ & Indices \\
Thm.\ 3 & Roots differ only by $T$ & Lemma 6 \\
Prop.\ 2 & Covariant: no difference & $P^2 = 1$ \\
\hline
\end{tabular*}}
\vspace{4pt}

% ======================================================================

\sloppy
\begin{thebibliography}{99}

\bibitem{NelsonLetter1}
S.~Nelson,
\textit{The Coordinate-Space Geometry of the Tetrad Frame},
Field Dynamics Research (2026), submitted.

\bibitem{NelsonLetter2}
S.~Nelson,
\textit{The Other Half of Gravity: Torsion and the Two-Chamber
Structure of the Tetrad Frame},
Field Dynamics Research (2026), submitted.

\bibitem{NelsonGramOrig}
S.~Nelson,
\textit{The Gram Matrix of the Tetrad Equation} (original note),
Field Dynamics Research (2026),
\url{papers/gram-matrix-of-tetrad-equation}.

\bibitem{NelsonScheme}
S.~Nelson,
\textit{The Tetrad Scheme and the Two Bounded Regions},
Field Dynamics Research (2026),
\url{proofs/tetrad-scheme-two-bounded-regions}.

\bibitem{NelsonForced}
S.~Nelson,
\textit{The Spatial Geometry of the Lorentz Tetrad},
Field Dynamics Research (2026),
\url{proofs/forced-tetrahedral-frame}.

\bibitem{NelsonHadamard}
S.~Nelson,
\textit{The Stellated Octahedron, the Cube and the Hadamard Matrix},
Field Dynamics Research (2026),
\url{proofs/hadamard-cube-frame}.

\bibitem{NelsonTorsion}
S.~Nelson,
\textit{The Torsion Sector of the Tetrad Connection as the Regular
Representation of $S_4$},
Field Dynamics Research (2026),
\url{proofs/torsion-regular-representation}.

\bibitem{NelsonChiral}
S.~Nelson,
\textit{Chiral Holonomy of the Closed Tetrad Frame},
Field Dynamics Research (2026),
\url{proofs/chiral-holonomy-closed-frame}.

\bibitem{LeviCivita1917}
T.~Levi-Civita,
\textit{Nozione di parallelismo in una variet\`{a} qualunque e
conseguente specificazione geometrica della curvatura Riemanniana},
Rend.\ Circ.\ Mat.\ Palermo \textbf{42} (1917) 173.

\bibitem{Cartan1922}
\'E.~Cartan,
\textit{Sur une g\'en\'eralisation de la notion de courbure de
Riemann et les espaces \`a torsion},
C.\ R.\ Acad.\ Sci.\ Paris \textbf{174} (1922) 593.

\bibitem{Hehl1976}
F.~W.~Hehl, P.~von der Heyde, G.~D.~Kerlick and J.~M.~Nester,
\textit{General relativity with spin and torsion: foundations and
prospects},
Rev.\ Mod.\ Phys.\ \textbf{48} (1976) 393.

\bibitem{Sylvester1867}
J.~J.~Sylvester,
\textit{Thoughts on inverse orthogonal matrices, simultaneous
sign-successions, and tessellated pavements in two or more colours},
Phil.\ Mag.\ \textbf{34} (1867) 461.

\bibitem{Kepler1619}
J.~Kepler,
\textit{Harmonices Mundi} (1619).

\bibitem{Coxeter1973}
H.~S.~M.~Coxeter,
\textit{Regular Polytopes},
3rd ed.\ (Dover, New York, 1973).

\end{thebibliography}

\end{document}
